\section{Conclusiones}

Los resultados obtenidos muestran que la codificación de fuente mediante el algoritmo de
Huffman logra una reducción significativa en la cantidad de bits necesarios para representar
el texto, respecto de un esquema de codificación de longitud constante: sobre el texto de
prueba utilizado, la codificación de Huffman requirió un \qty{42.12}{\percent} menos de bits
que el código ASCII extendido de 8 bits, alcanzando una eficiencia del \qty{99.44}{\percent}
respecto de la cota teórica de Shannon. Esto confirma la ventaja de asignar palabras de código
de longitud variable en función de la probabilidad de ocurrencia de cada símbolo, en lugar de
codificar todos los caracteres con la misma cantidad de bits.

Cabe destacar, sin embargo, que el texto utilizado no constituye, estrictamente, una fuente
discreta sin memoria: el modelo empleado supone que la aparición de cada caracter es
estadísticamente independiente de los caracteres anteriores, cuando en un texto real en inglés
existen fuertes correlaciones entre caracteres consecutivos (por ejemplo, la letra
\texttt{q} es seguida casi siempre por \texttt{u}). El modelo de fuente sin memoria utilizado
es, por lo tanto, una simplificación que ignora esta dependencia estadística; una codificación
que explote dicha memoria podría lograr una compresión aún mayor que la obtenida.

% vim: ts=4 sts=4 sw=4 et lbr
